\subsection{Нахождение обратной матрицы методом элементарных преобразований}

\begin{definition}
Матрица $B$ называется обратной к квадратной матрице $A$, если выполняются соотношения:
\[ BA = AB = E \]
где $E$ — единичная матрица. Обратная матрица обозначается $A^{-1}$.
\end{definition}

Для нахождения обратной матрицы используется метод элементарных преобразований строк. В основе метода лежат две теоремы.

\begin{theorem}
Квадратная матрица $A$ порядка $n$ может быть приведена к единичной матрице $E$ с помощью элементарных преобразований строк тогда и только тогда, когда она невырожденная (т.е. $\det A \neq 0$).
\end{theorem}

\begin{theorem}
Пусть $\det A \neq 0$. Если матрица $A$ превращается в единичную матрицу $E$ путём некоторого количества элементарных преобразований строк, то теми же самыми преобразованиями, применёнными в том же порядке, единичная матрица $E$ превращается в обратную матрицу $A^{-1}$.
\end{theorem}

\begin{tcolorbox}[title=Доказательство, breakable]
Согласно ранее доказанной теореме, каждое элементарное преобразование строк матрицы эквивалентно умножению этой матрицы слева на соответствующую матрицу элементарных преобразований.

Пусть последовательность из $k$ элементарных преобразований переводит матрицу $A$ в единичную матрицу $E$. Обозначим матрицы этих преобразований через $B_1, B_2, \dots, B_k$. Тогда результат последовательного применения этих преобразований к матрице $A$ можно записать в виде матричного равенства:
\[ B_k \cdot \dots \cdot B_2 \cdot B_1 \cdot A = E \]

Поскольку $\det A \neq 0$, существует обратная матрица $A^{-1}$. Умножим обе части полученного равенства на $A^{-1}$ справа:
\[ (B_k \cdot \dots \cdot B_2 \cdot B_1 \cdot A) \cdot A^{-1} = E \cdot A^{-1} \]
Пользуясь ассоциативностью умножения матриц и определением обратной матрицы ($A A^{-1} = E$), получаем:
\[ B_k \cdot \dots \cdot B_2 \cdot B_1 \cdot E = A^{-1} \]

Левая часть последнего равенства представляет собой результат применения той же самой последовательности элементарных преобразований к единичной матрице $E$. Правая часть равна $A^{-1}$. Таким образом, преобразования, переводящие $A$ в $E$, переводят $E$ в $A^{-1}$. Теорема доказана.
\end{tcolorbox}

\textbf{Алгоритм нахождения обратной матрицы (практический метод):}

\begin{enumerate}
    \item Составляется прямоугольная матрица размера $n \times 2n$ путем приписывания к матрице $A$ справа единичной матрицы $E$: $(A | E)$.
    \item С помощью элементарных преобразований только над строками (прямая и обратная прогонка метода Гаусса) левая часть этой матрицы (матрица $A$) приводится к единичному виду.
    \item В силу доказанной теоремы, в правой части на месте единичной матрицы получится искомая обратная матрица $A^{-1}$:
    \[ (A | E) \sim (E | A^{-1}) \]
    \item Если в процессе приведения матрицы $A$ к ступенчатому виду в левой части образуется чисто нулевая строка, это означает, что матрица $A$ вырождена ($\det A = 0$), и обратной матрицы для неё не существует.
\end{enumerate}

\begin{tcolorbox}[title=Источники, colback=gray!10]
\begin{itemize}
  \item Лекция №\,1, тема «Нахождение обратной матрицы методом элементарных преобразований», временной отрезок 36:00--45:00
  \item Лекция №\,2, тема «Теорема о вычислении обратной матрицы», временной отрезок 00:00--09:00
  \item PDF «Линейная алгебра\_compressed (1).pdf», раздел 3 «Вычисление обратной матрицы методом элементарных преобразований», стр. 5--7
  \item PDF «Вопросы», вопрос №\,3
\end{itemize}
\end{tcolorbox}
